\documentclass[12pt]{article}
\usepackage[utf8]{inputenc}
\usepackage[T1]{fontenc}
\usepackage{hyperref}
\title{Uzekwe Language Resources for K-12 Curriculum Development}
\author{Compiled by Claude Sonnet 5.5 (claude-sonnet-5-5) for Andr\'as Kornai}
\date{}
\begin{document}
\maketitle

\section{Introduction}
Uzekwe (also Ezekwe, Ezekwie, Zekwe; ISO 639-3 \texttt{eze}) is described as an Upper Cross River language of Cross River State, Nigeria (Ogoja region). Joshua Project reports about 19,000 people; Wikipedia cites an older figure of 5,000 speakers (1973). Very little material was found.

\section{Primary Educational Resources}
No primers, textbooks or school materials were found.

\section{Linguistic and Reference Materials}
\begin{itemize}
\item Zekwe language (Wikipedia): \url{https://en.wikipedia.org/wiki/Zekwe_language}
\end{itemize}

\section{Religious and Community Resources}
\begin{itemize}
\item Joshua Project people-group profile, Uzekwe in Nigeria (reports Bible translation as started): \url{https://joshuaproject.net/people_groups/15758/NI}
\item Joshua Project language page for eze: \url{https://joshuaproject.net/index.php/languages/eze}
\end{itemize}

\section{Digital and Audio Resources}
\begin{itemize}
\item Global Recordings Network listing for Ezekwe: \url{https://globalrecordings.net/en/language/eze}
\end{itemize}
The content of this listing was not inspected beyond its appearance in search results.

\section{Curriculum Development Context}
No evidence of school use was found. A curriculum would require new documentation work and community consultation.
\end{document}